\section{Appendix}
\label{sec:appendix}

\subsection{Repository-to-Paper Mapping}
\label{sec:repo_mapping}

Table~\ref{tab:repo_mapping} maps the main scientific concepts in the manuscript to the corresponding implementation modules. This appendix is intentionally concrete because the repository architecture is itself part of the reproducibility story.

\begin{table*}[h]
    \centering
    \scriptsize
    \begin{tabular}{>{\raggedright\arraybackslash}p{0.16\linewidth} >{\raggedright\arraybackslash}p{0.29\linewidth} >{\raggedright\arraybackslash}p{0.47\linewidth}}
        \toprule
        Paper concept & Main implementation modules & Role in the released code \\
        \midrule
        Frozen geometry prior & \texttt{training/factory.py}; \texttt{models/scene\_encoder/*}; \texttt{adapters/scene\_encoder/*} & Builds a frozen MapAnything or Depth Anything 3 encoder and exposes a uniform feature/reconstruction interface. \\
        Cross-view policy & \texttt{models/policy/attention\_policy\_network.py}; \texttt{geometry/pose\_ops.py} & Predicts a relative translation, transforms it to world coordinates, and converts it into a valid look-at pose. \\
        Batch preparation & \texttt{workflows/batch\_preparation\_use\_case.py}; \texttt{data/*} & Loads meshes, renders missing inputs when necessary, and selects the initial observation subset used by the policy. \\
        Candidate evaluation & \texttt{workflows/candidate\_evaluation\_use\_case.py}; \texttt{infrastructure/rendering/differentiable\_renderer.py} & Renders the predicted candidate view, fuses it with history, and constructs reconstruction inputs. \\
        Reconstruction-aware loss & \texttt{training/losses/reconstruction.py}; \texttt{chamfer\_regularizer.py}; \texttt{pose\_penalty.py} & Computes Chamfer-based reconstruction loss and geometric pose penalties. \\
        Training orchestration & \texttt{workflows/training\_step\_use\_case.py}; \texttt{training/lightning\_module.py} & Connects batch preparation, policy inference, candidate evaluation, optimization, logging, and test-time aggregation. \\
        Test protocol & \texttt{workflows/test\_evaluation\_use\_case.py}; \texttt{training/test\_metrics.py} & Computes GeomLoss, CD, DCD, and EMD for both the learned policy and a random-view baseline. \\
        \bottomrule
    \end{tabular}
    \caption{Mapping from paper concepts to code modules.}
    \label{tab:repo_mapping}
\end{table*}

\subsection{Recommended Finalization Checklist}
\label{sec:checklist}

Before turning this draft into a submission, we recommend the following concrete steps:
\begin{itemize}
    \item Run the default MapAnything + \texttt{point\_maps} configuration to convergence and populate Table~\ref{tab:main_results}.
    \item Run encoder and reconstruction-backend ablations to populate Table~\ref{tab:ablations}.
    \item Export qualitative results from the repository's observability pipeline and replace the placeholder figure in Fig.~\ref{fig:qualitative}.
    \item Decide whether the final paper should remain a frozen-prior NBV paper or become a VGGT-specific paper by wiring VGGT into the validated configuration and repeating the main experiments.
    \item Replace anonymous metadata, add acknowledgements, and expand the bibliography once the final baseline set is fixed.
\end{itemize}
